\section{Conclusion}
\label{sec:conclusion}

Power line communication remains an unusual but important communication technology. Its main advantage is clear: it can reuse conductors that already exist for power delivery. This can reduce installation effort and cost, especially in places where adding separate communication wiring would be expensive, inconvenient, or technically difficult.

At the same time, the limitations of PLC are equally clear. Power lines were not designed as communication channels. They are noisy, their impedance changes with connected loads and network conditions, and they require careful coupling and filtering to separate the communication signal from the power signal. In addition, PLC systems must deal with safety requirements and electromagnetic emission limits. Compared with more conventional communication media, such as dedicated data cables, optical fiber, or wireless systems, PLC often loses in speed, predictability, and reliability.

For these reasons, PLC is unlikely to become a universal or mainstream communication solution. It is not the best choice when a cleaner, faster, or more secure communication medium is easily available. However, the examples discussed in this paper show that PLC can be very useful in carefully selected cases. Home networking, smart metering, utility control, and industrial feedback over existing power wiring all demonstrate the same basic idea: when the power connection already reaches the required device, reusing it for communication can be a practical engineering compromise.

Therefore, PLC should not be understood as an obsolete technology or as a universal replacement for other communication systems. It is better understood as a niche technology with persistent value. Its challenges prevent it from becoming dominant, but it will remain useful as long as we use wires to transmit energy.